% ============================================================
% CAPITOLO - METODO DI LAVORO PER IL BASSO NUMERATO
% (Versione espansa di 99_basso_dato.tex)
% ============================================================
\chapter{Come affrontare il basso numerato}
\label{chap:metodo_lavoro}

L'armonizzazione di un basso numerato è la prova fondamentale di ogni esame di armonia. In questo capitolo presentiamo un \textbf{metodo di lavoro sistematico}, una procedura in fasi sequenziali pensata per trasformare un esercizio potenzialmente disorientante in un percorso logico e controllato.

Il basso numerato (o \textbf{basso figurato}) è un basso con le cifre che indicano gli intervalli da costruire sopra di esso. Rispetto al basso senza numeri (Capitolo \ref{chap:dal_numerato_al_senza_numeri}), la numerica ci fornisce già la ``numerica implicita'': il nostro compito è tradurre quei numeri in accordi reali, rispettando la condotta delle voci.

\section{Lettura della numerica: convenzioni fondamentali}
\label{sec:lettura_numerica}

\begin{definizione}[La numerica del basso figurato]
	\label{def:numerica}
	La \textbf{numerica} (o \textbf{cifratura}) è un sistema di numeri scritti sotto le note del basso che indicano gli intervalli delle voci superiori rispetto al basso stesso. Le convenzioni principali sono:
	\begin{itemize}
		\item \textbf{Nessun numero} (o \textbf{5}): accordo in stato fondamentale (quinta e terza sopra il basso).
		\item \textbf{6} (o 6/3): primo rivolto (sesta e terza sopra il basso).
		\item \textbf{6/4}: secondo rivolto (sesta e quarta sopra il basso) --- usare solo nei contesti ammessi (cfr. Capitolo \ref{chap:quarta_sesta}).
		\item \textbf{7}: accordo di settima allo stato fondamentale.
		\item \textbf{6/5}: primo rivolto di settima (sesta, quinta e terza sopra il basso).
		\item \textbf{4/3}: secondo rivolto di settima.
		\item \textbf{2} (o 4/2): terzo rivolto di settima (la settima è al basso).
		\item Un \textbf{numero preceduto da \#, \flat{} o $\natural$}: l'alterazione si applica all'intervallo indicato.
		\item Un \textbf{\# o + da solo} sotto una nota: alzare di un semitono la terza dell'accordo (comune in Do minore per alzare la sensibile).
	\end{itemize}
\end{definizione}

\begin{osservazione}[La numerica è relativa al basso, non alla scala]
	\label{oss:numerica_relativa}
	Gli intervalli della numerica sono \textit{relativi al basso reale}, non al grado della scala. Se il basso è Re e la numerica è 6, si deve costruire una sesta sopra Re (Si), indipendentemente dalla tonalità. Le alterazioni esplicite nella numerica hanno sempre la precedenza sulle note ``naturali'' della scala.
\end{osservazione}

\section{Fase 1 — Analisi preliminare}
\label{sec:fase1_numerato}

Anche con la numerica disponibile, la prima operazione non è scrivere accordi, ma \textbf{comprendere il terreno}.

\begin{itemize}
	\item \textbf{Determinare la tonalità:} Osservare l'armatura di chiave e verificare la nota iniziale e finale. Prestare attenzione alle alterazioni cromatiche presenti nel basso: spesso indicano la sensibile del modo minore (\ref{chap:tonicizzazione}) o modulazioni (\ref{chap:modulazione}).
	\item \textbf{Individuare il modo:} Cercare la presenza della sensibile (VII grado alzato nel minore) e le note tipiche del minore armonico.
	\item \textbf{Leggere la numerica:} I numeri sono \textit{suggerimenti} preziosi. Indicano la struttura degli accordi desiderata dal compositore/esaminatore, ma vanno sempre interpretati nel contesto della tonalità e della condotta delle voci.
	\item \textbf{Riconoscere le note non armoniche:} Identificare eventuali note di passaggio, note di volta o bordoni nel basso stesso (\ref{chap:noteestranee}). Questi movimenti non portano accordi propri.
\end{itemize}

\section{Fase 2 — Segmentazione fraseologica}
\label{sec:fase2_numerato}

Dividere il basso in unità di senso compiuto permette di pianificare i punti di riposo.

\begin{itemize}
	\item \textbf{Individuare le cadenze:} Cercare i movimenti tipici (V-I, IV-I, formule cadenzali) che segnano la fine di una frase o di un periodo. Spesso la numerica stessa rivela la cadenza: una successione di numeri \texttt{<4> <3>} indica un ritardo 4-3 cadenzale (\ref{chap:ritardi}); \texttt{<7>} sulla dominante indica la settima di dominante in cadenza.
	\item \textbf{Struttura antecedente-conseguente:} Verificare se il basso segue la tipica struttura di 8 misure (\ref{chap:struttura_frase}). Pianificare in anticipo la cadenza di metà (aperta) e quella finale (perfetta).
\end{itemize}

\begin{regola}[Pianificazione cadenzale nel basso numerato]
	\label{reg:pianificazione_cadenzale_numerato}
	\begin{itemize}
		\item \textbf{Metà (mis. 4 di 8):} Deve presentare una cadenza sospesa (V) o imperfetta (I non fondamentale o soprano non in ottava).
		\item \textbf{Finale (mis. 8):} Deve presentare una cadenza perfetta (V-I in stato fondamentale, soprano all'ottava o alla quinta).
		\item Se la numerica non specifica risoluzione della dominante: usare la forma completa (V7 completa) o incompleta secondo il contesto (\ref{def:v7_completa_incompleta}).
	\end{itemize}
\end{regola}

\section{Fase 3 — Classificazione di ogni nota di basso}
\label{sec:fase3_numerato}

Stabilire la funzione e lo stato di ogni accordo, \textit{leggendo la numerica e interpretandola nel contesto}.

\begin{itemize}
	\item \textbf{Assegnazione degli accordi:} Decidere quale accordo assegnare a ogni nota in base alla funzione armonica (\ref{chap:funzioni_armoniche}) e alla numerica disponibile. Verificare che l'accordo indicato dalla numerica sia coerente con la tonalità.
	\item \textbf{Tonicizzazioni e cromatismi:} Identificare se un'alterazione nella numerica indica una dominante secondaria (\ref{chap:tonicizzazione}) o una modulazione (\ref{chap:modulazione}). Un \# sulla terza dell'accordo spesso indica la sensibile di una tonalità locale.
	\item \textbf{Scelta dello stato:} La numerica specifica lo stato, ma va verificata la coerenza con:
		\begin{enumerate}
			\item La posizione metrica (stati fondamentali sui tempi forti di norma).
			\item L'uso corretto del secondo rivolto (\ref{chap:quarta_sesta}).
			\item Le regole di risoluzione delle dissonanze (\ref{chap:teoria_dissonanza}).
		\end{enumerate}
\end{itemize}

\section{Fase 4 — Costruzione delle voci superiori}
\label{sec:fase4_numerato}

Una volta stabilita l'ossatura armonica, si procede alla scrittura SATB.

\begin{regola}[Priorità nella scrittura delle voci superiori]
	\label{reg:priorita_scrittura_numerato}
	\begin{enumerate}
		\item \textbf{Vettori obbligati:} Risolvere sempre correttamente la sensibile (sale alla tonica) e la settima (scende di grado). Questi movimenti hanno \textit{priorità assoluta} sugli altri.
		\item \textbf{Note comuni:} Tenere le note comuni tra accordi successivi nella stessa voce.
		\item \textbf{Grado congiunto:} Muovere le altre voci preferibilmente per grado congiunto.
		\item \textbf{Moto contrario:} Preferire il moto contrario tra basso e soprano per evitare quinte e ottave nascoste.
	\end{enumerate}
\end{regola}

\begin{itemize}
	\item \textbf{Scelta dei raddoppi:} Raddoppiare di norma la fondamentale. \textbf{Mai} raddoppiare la sensibile o la settima.
	\item \textbf{Posizione al soprano:} Scegliere la posizione (ottava, terza, quinta) cercando di creare una linea melodica cantabile. La nota indicata dalla numerica deve comparire nelle voci superiori, ma non necessariamente al soprano.
\end{itemize}

\subsection{La gestione delle dissonanze nel basso numerato}

Quando la numerica indica una dissonanza (settima, ritardo, sesta eccedente):

\begin{regola}[Trattamento delle dissonanze nella realizzazione]
	\label{reg:dissonanze_realizzazione}
	\begin{enumerate}
		\item \textbf{Identificare la nota dissonante}: qual è la nota che crea la dissonanza rispetto al basso?
		\item \textbf{Verificare la preparazione}: la nota dissonante era consonante nell'accordo precedente nella stessa voce? (Obbligatorio per II7, IV7, ritardi; non necessario per V7.)
		\item \textbf{Pianificare la risoluzione}: la nota dissonante deve scendere di grado nell'accordo successivo. Quale nota dell'accordo successivo è la sua destinazione?
		\item \textbf{Evitare conflitti di risoluzione}: se due note devono risolvere contemporaneamente, assicurarsi che le loro risoluzioni non producano parallelismi o ottave.
	\end{enumerate}
\end{regola}

\section{Fase 5 — Revisione dell'elaborato}
\label{sec:fase5_numerato}

L'ultima fase è la più critica: il controllo degli errori.

\begin{regola}[Protocollo di revisione del basso numerato]
	\label{reg:revisione_numerato}
	Utilizzare sistematicamente la checklist degli errori (\ref{reg:checklist_errori}). Suggerire di rileggere ogni coppia di voci separatamente (S+B, A+T, S+A, T+B) prima di valutare l'insieme.
	
	Attenzione particolare a:
	\begin{itemize}
		\item La sensibile al soprano ha sempre risolto verso l'alto?
		\item Ogni settima si è preparata (dove necessario) e risolta in discesa?
		\item Ogni ritardo indicato nella numerica è stato realizzato nel tempo forte corretto?
		\item Il secondo rivolto è stato usato solo nei contesti ammessi (\ref{chap:quarta_sesta})?
		\item I raddoppi rispettano il principio (mai sensibile, mai settima)?
	\end{itemize}
\end{regola}

Il segreto di un buon basso non è solo l'assenza di errori, ma la \textbf{coerenza del discorso musicale} che riusciamo a costruire sopra le note date. La numerica ci vincola nella scelta degli accordi, ma ci lascia liberi nella condotta delle voci e nella qualità melodica del soprano.

\section{Esempi di realizzazione commentati}
\label{sec:esempi_realizzazione}

\subsection{Basso numerato con settima di dominante e ritardo}

\musica{
	\new Staff {
		\clef bass \key do \major \time 4/4
		\relative do, {
			do2 fa | sol2 do
		}
	}
}
Numerica: \texttt{s2 s | <7>2 <4 3>}

\textbf{Lettura:}
\begin{itemize}
	\item Battuta 1: Do (I, 5) -- Fa (IV, 5). Nessun numero = stati fondamentali.
	\item Battuta 2: Sol con 7 (V7) -- Do con 4-3 (I con ritardo 4-3).
\end{itemize}

\textbf{Decisioni di condotta:}
\begin{itemize}
	\item V7: basso Sol; sensibile (Si) al tenore o soprano; settima (Fa) deve scendere a Mi.
	\item Ritardo 4-3: il Fa (quarta sopra Do) è già presente nel V7 come settima. Viene tenuto e risolve su Mi (terza di I). Questo è il ritardo preparato dalla settima.
	\item Il soprano può avere il Do (ritardo 4, poi Mi, risoluzione 3) oppure un'altra voce porta il ritardo.
\end{itemize}

\musica{
	\partilate
	{ \key do \major \time 4/4 }
	{ mi'2 fa' | \freccia fa'4 \glissando mi'4 mi'2 }
	{ do'2 do' | re'2 do'2 }
	{ sol2 la | si2 sol2 }
	{ do2 fa, | sol,2 do2 }
	\lyricmode { "I"2 "IV" | "V7"2 "I"2 }
	\figuremode { s2 s | <7>2 <4 3>2 }
}

\subsection{Basso numerato con modulazione implicita}

\musica{
	\new Staff {
		\clef bass \key re \major \time 4/4
		\relative re, {
			re2 mi | fad2 sol | la2 si | dod1
		}
	}
}
Numerica: \texttt{s2 s | <6>2 s | s2 s | <6>1}

\textbf{Analisi:}
\begin{itemize}
	\item Tonalità: Re maggiore (\Fadiesis{}, \Dodiesis{}).
	\item Battuta 4: \Dodiesis{} con 6 $\rightarrow$ accordo di \Dodiesis{} con la terza al basso. \Dodiesis{} è la sensibile di Re. L'accordo sarà VII°6 di Re maggiore $\rightarrow$ rimaniamo in Re maggiore.
	\item Il brano sale per gradi congiunti (I-II-III-IV-V-VI-VII): Regola dell'Ottava completa ascendente in Re maggiore.
\end{itemize}

\musica{
	\partilate
	{ \key re \major \time 4/4 }
	{ fad'2 fad' | fad' sol' | la'2 si' | do''1 }
	{ re'2 re' | re' re' | fad'2 sol' | la'1 }
	{ la2 si | la si | do'2 re' | fad'1 }
	{ re,2 mi, | fad, sol, | la,2 si, | dod1 }
	\lyricmode { "I"2 "V" | "I6" "II6" | "V"2 "III" | "VII°6"1 }
	\figuremode { s2 s | <6>2 <6> | s2 s | <6>1 }
}
